\begin{abstract}
Coding agents increasingly operate across sessions, tools, and handoffs, yet a
strong fresh model can often recover a project directly from current source and
tests. This makes a useful evaluation question narrower than whether more
history is always better: when does project memory improve a dependent task, and
when should an agent ignore it? We present a preregistered evaluation artifact
for freshness-aware multi-session project memory, MemorixBench-Transfer; this
is a protocol and reproducibility study, not a completed efficacy study. The
design compares a
capable no-memory control with bounded transcript replay, equal-evidence
retrieval adapters, and a separately reported native \memorix{} MCP surface.
It records source freshness, provenance, task-level correctness, context cost,
and observable actions that conflict with frozen stale claims. Its confirmatory
protocol specifies isolated workers, private controller-side oracles,
independent case admission, and separate worker/relay/runtime evidence chains;
the current public cohorts do not claim those deployment gates. Repeated runs
estimate within-case stability rather than inflate sample size. On a frozen 12-fixture public Qwen
cohort with 144 valid runs, canonical \memorix{} achieves a $+2.8$ point
case-clustered success difference over no memory, with a 95\% interval spanning
$-8.3$ to $+16.7$ points. A separately frozen 72-run DeepSeek replication has
100\% task success under both conditions and added context cost for canonical
\memorix{}. These public results are descriptive rather than confirmatory:
the fixtures cannot establish pretraining novelty or private-oracle integrity,
and the KVM-backed execution cohort remains under construction.
\end{abstract}
